% 论文⑤中文主本。英文对照版后行。
% 编译：tectonic paper_zh.tex（离线：tectonic --only-cached paper_zh.tex）
\documentclass[10pt,fontset=fandol]{ctexart}
\usepackage[margin=1in]{geometry}
\usepackage{amsmath}
\usepackage{amssymb}
\usepackage{booktabs}
\usepackage{tabularx}
\usepackage{longtable}
\usepackage[numbers]{natbib}
\usepackage{tikz}
\usetikzlibrary{arrows.meta,positioning}
\usepackage[hidelinks]{hyperref}

\newcommand{\en}[1]{（#1）}

\title{\textbf{四个概念，四种类型：\\智能、人类智能、人工智能、通用人工智能的\\层次、边界与转译协议}}
\author{余玉桂\\ 余玉桂战略研究室\\ \texttt{they4897@gmail.com}}
\date{}

\begin{document}
\maketitle

\begin{abstract}
\noindent 智能、人类智能、人工智能、通用人工智能——四个词在文献、政策与产业话语中每天使用，但它们之间的层次与边界至今没有权威版本：仅一部编纂就收录了约七十条智能定义；通用人工智能的定义在 2023 至 2026 年间进入制度化竞争期，产业界则通过系统性换词，避免被自己立下的定义反噬。本文提出双重主张。其一（类型论）：四个词不是一条阶梯上的四级，而是四种不同类型的概念——智能是能力概念、人类智能是样本概念、人工智能是造物概念、通用人工智能是参照系概念；每一类的“定义完成”形态不同，流行的“专用—通用—超级”阶梯是把四类概念串在一条线上的叙事装置。其二（分层论）：每个概念内部须区分本体层（它是什么）、认识层（什么证据支持某系统具有它）与方法层（用什么程序检测它）；定义之争长期难解的根源是三种非法滑动——操作主义（测量冒充定义）、部分谬误（机制冒充能力定位）、判据升格（工程判据冒充概念内容）。本文的核心贡献是一份转译协议：工程观察（电路、特征、梯度、基准得分）到概念命题的转译表，每行附生效条件，附非法转译目录；据该协议，“泛化在权重中”类主张不是错误而是越权——能力谓词的归属层级是系统与环境范围的组合，关系属性本无所谓位置。文末给出四词在 AGI 人本拟合论中的最终位置、证伪条款与中文语境对接。
\end{abstract}

\medskip
\noindent\textbf{关键词：} 智能定义；通用人工智能；构造效度；操作主义；部分谬误；机制解释性；概念分层；AGI 人本拟合论

\section{引言}

这四个词在人工智能话语中被日日使用，其间的层次关系却从未有过共识。“什么是智能”的问题至少从 1923 年起就没有共识\cite{boring1923intelligence}；Legg 与 Hutter 编纂过的非正式定义约七十条\cite{legg2007definitions}——定义的数量在一个世纪里增加而非收敛。“什么是通用人工智能”（AGI）使同一困境雪上加霜：2023 至 2026 年，AGI 的定义进入制度化竞争期——性能分级矩阵\cite{morris2023levels}、心理测量模型定义\cite{hendrycks2025definition}、经济条款定义\cite{openai2018charter}并存而互斥；与此同时，产业与监管话语系统性换词（前沿模型、超级智能、通用目的人工智能），没有任何主要机构愿意受自己五年前的定义约束、对“AGI 是否已经达成”作出裁决。四词之间唯一深入人心的“层次图”——专用人工智能到 AGI 再到超级智能的三级阶梯——是大众惯例而非学术分类，而且，本文将论证，它本身就是概念混淆的载体。

两类日常话语让这种无共识的状态变得刺眼。第一类来自定义侧：大量“定义”实际上是测量方法。“智能就是测验所测的东西”\cite{boring1923intelligence}是这一形态的原型；它的当代变体包括把一个测量泛函当作智能的定义\cite{legg2007universal}、把一套评测协议当作通用性的定义\cite{morris2023levels,chollet2019measure}、把一份经济条款当作 AGI 的定义\cite{openai2018charter}。第二类来自工程侧：机制解释性研究在权重、电路与梯度中观察到结构，随后在传播中升格为“泛化、智能就在其中”的断言。本文对这两类话语给出的诊断相同：它们都不是错误命题，而是\emph{越权命题}——把在一层之内合法的主张，未经条件检验就转译到另一层。而“让人难以信服”的结构性原因也在于此：方法冒充定义的命题不可反驳——测试失败永远可以归因于方法（基准污染、分布不匹配、任务集过窄），定义本身从不受伤；一个永不承担失败风险的命题不是认识论主张，是信条。

本文做三件事。第一，把四个词整理成四种概念类型，并说明为什么每一类的“定义完成”形态不同（第~\ref{sec:types} 节）。第二，给每个概念建立本体层、认识层、方法层的三层结构，把现有定义之争诊断为三种非法滑动，并以一份逐字引语的主张清单为证据底座（第~\ref{sec:layers} 节）。第三，给出一份转译协议——工程观察到概念命题的转译表与非法转译目录（第~\ref{sec:protocol} 节）；这是本文的核心贡献：它不禁止跨层转译，而是要求每次转译出示生效条件。第~\ref{sec:program} 节把四词安放进 AGI 人本拟合论，第~\ref{sec:falsify} 节给出证伪条款，附录~A 给出主张清单，附录~B 给出中文语境对接。

本文与同一研究纲领前三篇论文分工明确：理论论文论证了通用性是相对于人类实践开放集的关系属性\cite{yu2026theory}，本文把这一结论用作参照系概念的现成理论形式；概念解绑论文区分了人工智能的工程目标与机器之心\cite{yu2026twogoals}，本文把它的“功能智能”用作造物概念的判据；方法论论文的披露审计\cite{yu2026methodology}在本文中用作认识层工具的范例。本文的增量贡献是四词的类型学、三层结构与转译协议——经文献检索未见先例（中英文文献中均无，见第~\ref{sec:survey}、\ref{sec:cn} 节的空白声明）。

\section{定义现状：谱系与漂移}\label{sec:survey}

\subsection{智能：无共识的世纪}

现代讨论仍在引用的两极，早在 1904 至 1923 年间就已成形。一极是心理测量学：Spearman 从“几乎所有认知测验彼此正相关”的统计事实出发提出一般因素 $g$\cite{spearman1904general}，由此开启了一条在人类样本内工作得相当好的测量传统。另一极是操作主义立场：Boring 直言“智能就是智力测验所测的东西”\cite{boring1923intelligence}——定义即测量操作。Bridgman 在物理学中提出“概念即一套操作”\cite{bridgman1927logic}后，操作主义在物理科学中迅速退潮，却在智能定义史上反复出现（第~\ref{sec:layers} 节）。

形式化传统的代表是 Legg 与 Hutter：“智能衡量的是一个智能体在广泛环境范围内达成目标的能力”\cite{legg2007universal}，并以对一切可计算环境的加权期望奖励给出形式定义。这一路径把智能定义为一种与载体无关的能力度量，人类智能在该框架内降格为尺度上的一点。与之相对，Chollet 把智能定义为“在给定任务范围、先验与经验条件下获取新技能的效率”\cite{chollet2019measure}，并断言技能只是智能过程的固定产物——不存在这样的任务：仅凭系统在该任务上的技能就足以证明智能。两个定义互相冲突：前者会把一个“买来大量技能”的系统判为高智能，后者判其为高级技能库。一个世纪过去，分歧的形态没有收敛，只是换了词汇。

\subsection{人类智能：样本内的共识}

对“人类智能”，心理学反而有一份在实务层面足以使用的共识。Cattell--Horn--Carroll 三层模型把能力结构整理为：顶层 $g$、中层约十个广域能力（流体推理、晶体知识、加工速度、工作记忆等）、底层数十种窄能力\cite{mcgrew2023chc}；当代主流智力测验均按该模型设计。两份共识文件给出了引用最广的文字定义：美国心理学会任务组报告\cite{neisser1996intelligence}与五十二名研究者联署的《关于智能的主流科学》\cite{gottfredson1997mainstream}。

这份共识的边界常被忽略：它只覆盖智人样本。Legg 与 Hutter 编纂的七十条定义中多数来自心理学家、以人类为样本——“许多智能定义实为人类智能定义”并非修辞，而是可以直接从这部汇编中核实的事实\cite{legg2007definitions}。Hern\'andez-Orallo 的通用心理测量学系统论证了：以人类为标定参照系的测验，外推到机器与动物时便会失效\cite{hernandezorallo2017measure}；Chollet 则给出一个重要的反转：人类中心的参照系“不仅合法而且必要”——因为技术为人服务——但必须明确作出声明，把先验与开发者注入计入度量\cite{chollet2019measure}。也就是说，学界已经承认“智能”一词在人类中心与载体无关两种立场间摇摆，但从未把这一摇摆本身整理成概念类型的区分。这正是本文要做的事。

\subsection{人工智能：不定义智能，只定义制造}

人工智能的奠基文本有意绕开了智能的定义。达特茅斯提案的纲领句是：学习的每一方面或智能的任何特征，原则上都可以精确描述，从而使机器能够模拟它\cite{mccarthy1955dartmouth}——智能是不加定义的初始词项，被定义的是“制造”。McCarthy 本人的正式表述延续这一结构：“人工智能是制造智能机器、尤其智能计算机程序的科学工程”\cite{mccarthy2007whatis}。教科书传统把定义空间整理为“像人一样思维/理性地思维、像人一样行动/理性地行动”的四方矩阵，并选定“理性地行动”一支\cite{russell2021aima}；图灵测试一脉则把人设为不可分辨的对手\cite{turing1950computing}。

这一“制造优先”结构有一个著名的后门：由于智能本身未被定义，“人工智能”的边界只能随工程成就移动——Tesler 定理“智能就是机器尚未做到之事”\cite{tesler1997effect}刻画的正是这个否定性、移动的类别。它解释了为什么 McCarthy 的定义能兼容任何具体技术，也预示了本文第~\ref{sec:layers} 节要诊断的第三种滑动：当工程判据（基准、产品、商业成功）足够强势时，它们便会在不知不觉间被当作智能本身的定义。

\subsection{通用人工智能：制度化竞争期}

AGI 的定义谱系在 2023 年前是学术小传统：Legg 与 Hutter 的定义与“智能”同一条\cite{legg2007universal}；Goertzel 把它表述为“在复杂环境中达成复杂目标的能力”\cite{goertzel2014agi}。2023 年起进入制度化竞争：GPT-4 的“火花”论文以“广泛认知能力”作隐含定义\cite{bubeck2023sparks}；DeepMind 的分级框架给出“性能$\times$通用性”矩阵，把“通用”操作化为“达到目标性能阈值的任务范围”，并将当前前沿模型评为“新兴 AGI”\cite{morris2023levels}；OpenAI 章程给出经济定义——“在大多数具有经济价值的工作上超越人类的高度自主系统”\cite{openai2018charter}；2025 年，三十三位作者联署的定义以 Cattell--Horn--Carroll 心理测量模型为骨架，把 AGI 定义为“匹配受良好教育的成年人的认知多面性与熟练度”\cite{hendrycks2025definition}。批评侧，有论证指出 AGI 话语的诸定义“拥抱一系列互不相容的价值与假设”\cite{blili2024unsocial,blili2025stop}；测量侧，新提案继续涌现\cite{xu2024meant,dhar2025generality,gignac2024defining}。

同期发生了两重用词漂移：产品与监管语言从 AGI 转向“前沿模型”与“通用目的人工智能”（后者以训练算力阈值推定系统性风险）；实验室叙事从 AGI 转向“超级智能”。定义分歧最终以换词的方式被搁置，而非得到解决。

\subsection{中文语境}\label{sec:cn}

中文语境为四词提供了两个独有的观察角度。

第一条是术语分工史。心理学自 1924 年陆志韦\cite{zhang2020psychometrics}首次修订比纳量表起，固定使用“智力”（智力测验、智商量表）；“人工智能”译名于二十世纪七十年代末传入、八十年代初在学界正式启用，固定使用“智能”。“智力—智能”的学科分工格局沿用至今，辅以科普层面的“智力（基础层）—智能（应用层）—智慧（价值层）”三层递进——但这一三层框架从未得到严格的概念分析。换言之，中文语境对“智能”的分层有长期的直觉，缺的是学理化的层次图。

第二条是政策用词的漂移。2017 年《新一代人工智能发展规划》仅在基础理论布局中提及“通用人工智能”\cite{state2017plan}；2023 年北京的地方政策在操作层面把通用人工智能等同于“通用大模型技术体系”\cite{beijing2023measures}——与学理定义（如朱松纯团队的“无限任务、自主定义任务、价值驱动”三特征与通智测试体系\cite{zhu2025tongtest}）之间存在显见的落差；以 2025 年政府工作报告为代表的中央文件则完全绕开“通用人工智能”一词，改用“人工智能+”“具身智能”“智能经济”等政策热词\cite{gov2025report}。国内学界的定义性工作自成谱系：张钹的三代人工智能论给出“专用—通用”的历史分期\cite{zhang2020third}；钟义信的机制主义理论给出信息—知识—智能转换的中文原创定义\cite{zhong2018mechanism}；哲学侧，刘永谋论证通用人工智能陷入“唯名论困境”——“存在着具体的 AGI 项目，却没有关于 AGI 的一致的定义”\cite{liu2024complex}；赵汀阳与殷杰则从主体性维度切入\cite{zhao2019selfconscious}。经文献检索，中文文献中未见以四词层次与边界为题的系统分析（附录~B 给出对话者定位与空白声明）。

\subsection{小结：层次的空缺}

把谱系合起来看：“智能”有两大互斥定义传统；“人类智能”有样本内共识但常被冒充一般定义；“人工智能”有意不定义智能只定义制造；“通用人工智能”处于定义的制度化竞争期并以换词回避裁决。四词层次关系的权威版本不存在；通用性有四种互不兼容的操作化（任务范围、环境开放性、技能获取效率、自主性等级）；人类智能在 AI 定义中扮演四种角色（参照系、对手、经济计量单位、被批判的标尺）。层次的空缺不是文献不足，而是缺少一个能把这四个词各自属于哪种概念类型说清楚的框架。第~\ref{sec:types} 节给出这个框架。

\section{论题一：四个词，四种概念类型}\label{sec:types}

\subsection{判据}

本文用以分类的判据是：一个概念的类型由其成功判据的形态决定——即“定义问题已被回答”对它意味着什么。四词在该判据下落入四类（图~\ref{fig:matrix}）。

\begin{figure}[t]
\centering
\begin{tikzpicture}[font=\small,
  box/.style={draw, rounded corners=2pt, minimum width=3.05cm, minimum height=1.05cm, align=center},
  slide/.style={-{Stealth[length=2.2mm]}, dashed, thick, red!70!black}]
% 三层标签
\node[anchor=east, font=\bfseries] (ont) at (-0.25,4.2) {本体层};
\node[anchor=east, font=\bfseries] (epi) at (-0.25,2.8) {认识层};
\node[anchor=east, font=\bfseries] (met) at (-0.25,1.4) {方法层};
% 列标签
\node[font=\bfseries] at (1.9,5.35) {智能};
\node[font=\bfseries] at (5.15,5.35) {人类智能};
\node[font=\bfseries] at (8.4,5.35) {人工智能};
\node[font=\bfseries] at (11.65,5.35) {通用人工智能};
% 智能列
\node[box] (i1) at (1.9,4.2) {能力谓词\\系统 $\times$ 环境范围};
\node[box] (i2) at (1.9,2.8) {跨分布的证据\\构造效度};
\node[box] (i3) at (1.9,1.4) {测验 / 泛函 / 协议};
% 人类智能列
\node[box] (h1) at (5.15,4.2) {样本的测量结构\\（$g$、广域能力）};
\node[box] (h2) at (5.15,2.8) {常模参照\\信度、效度};
\node[box] (h3) at (5.15,1.4) {标准化测验};
% AI 列
\node[box] (a1) at (8.4,4.2) {造物：按构造与\\功能判据定义};
\node[box] (a2) at (8.4,2.8) {工程证据\\公开、密集};
\node[box] (a3) at (8.4,1.4) {规格、基准、产品判据};
% AGI 列
\node[box] (g1) at (11.65,4.2) {参照系：通用性相对于\\人类实践开放集};
\node[box] (g2) at (11.65,2.8) {实践抽样的明确声明\\（不可闭合）};
\node[box] (g3) at (11.65,1.4) {分级矩阵 / 审计\\（皆为替代测量）};
% 三种滑动箭头（正交走线：沿图边与列间空隙，不穿格子）
\draw[slide, rounded corners=4pt] (i3.west) -- (0.05,1.4) -- (0.05,4.2) -- (i1.west);
\node[font=\footnotesize, red!70!black, anchor=east] at (-0.25,3.5) {操作主义};
\draw[slide, rounded corners=4pt] (a3.south) -- (8.4,0.3) -- (0.2,0.3) -- (0.2,3.675) -- (i1.south west);
\node[font=\footnotesize, red!70!black] at (4.5,0.57) {部分谬误：机制 $\to$ 能力定位};
\draw[slide, rounded corners=4pt] (a3.east) -- (10.02,1.4) -- (10.02,4.2) -- (g1.west);
\node[font=\footnotesize, red!70!black] at (11.65,4.95) {判据升格};
\end{tikzpicture}
\caption{四个概念$\times$三层结构。列即概念类型（能力、样本、造物、参照系），行即层次（本体层、认识层、方法层）。红色虚线标出第~\ref{sec:layers} 节诊断的三种非法滑动：方法层主张被升格为本体层定义（操作主义）；工程观察被升格为能力的定位（部分谬误）；造物判据被升格为能力定义与参照系内容（判据升格）。}
\label{fig:matrix}
\end{figure}

\subsection{智能：能力概念}

“智能”谓词的归属层级是\emph{系统与环境的组合}：说一个系统有智能，总是相对于一个环境范围而言。Legg 与 Hutter 的形式定义把这一点表述为数学——智能是环境空间上的期望性能\cite{legg2007universal}；心理测量学的做法是把它限制在人类所遇任务总体之内。两种进路的共同点：智能的成功判据是\emph{度量}形态——定义问题得到回答，意味着存在一个带构造效度声明的度量。作为一个能力概念，智能不回答“系统由什么构成”，只回答“系统在什么范围内能做到什么”。

\subsection{人类智能：样本概念}

“人类智能”不是另一个能力，而是\emph{一个被测群体的测量结构}。它的成功判据是常模形态：定义问题得到回答，意味着对一个有定义的人群存在标准化工具、常模与信效度证据\cite{spearman1904general,mcgrew2023chc,neisser1996intelligence}。样本概念与能力概念的关系是投影而非子集：人类智能是智能概念在智人样本上的投影，带有人群边界——这解释了第~\ref{sec:survey} 节的两件事：为什么心理测量共识不能外推到机器（投影换了样本）\cite{hernandezorallo2017measure}，以及为什么“许多智能定义实为人类智能定义”（把投影当成了原象）。

\subsection{人工智能：造物概念}

“人工智能”是一种造物：按构造与功能判据定义的工程物类别。它的成功判据是工程形态——功能规格达成、可公开检验的性能——不要求“拥有智能”作为成立条件。搜索引擎是人工智能，无论有没有人打算把“智能”谓词授予它。本文纲领的概念解绑论文已给出该层的内容：人工智能的工程目标是服务于人类实践的功能智能，其成功判据公开、密集、持续改进\cite{yu2026twogoals}。关键在“造物”二字的推论：造物概念的判据（好不好用）\emph{不}定义能力概念的内容（什么是智能）——这一推论在第~\ref{sec:layers} 节成为第三种滑动的诊断依据。

\subsection{通用人工智能：参照系概念}

“通用人工智能”的类型最反直觉：它不是一个更大的能力，而是一个\emph{参照系}。理论论文已论证其形式：通用性不是自足的内在属性，而是关系属性——一个系统称得上“通用”，只相对于它所参照的人类实践开放集成立；“通用”参照实践的开放集，“高”参照人类教师的绩效，两者都是需求一方作出的规定\cite{yu2026theory}。参照系概念的成功判据不可闭合：它不能被一份规格“定义完成”，只能以明确声明的方式给出——声明参照哪个实践集、以什么抽样近似它。这正是定义分歧无法终结的结构性原因：每一份 AGI 定义实际上都在声明一个参照系（CHC 能力清单\cite{hendrycks2025definition}、经济价值工作集\cite{openai2018charter}、任务广度阈值\cite{morris2023levels}），而参照系之间的分歧不是经验分歧，没有实验能裁决。

\subsection{阶梯批判与四关系图}

流行的三级阶梯（专用人工智能$\to$AGI$\to$超级智能）把四个概念类型串在一条单变量线上，仿佛存在一个“智能量”沿阶梯增长。按类型论，正确的图景是四种概念加四种关系：人类智能是智能在人类样本上的投影（测量关系）；人工智能是实现智能候选者的造物类别（实现关系）；对人工智能的性能评估调用人类智能作为基准（参照关系）；通用人工智能是对参照系本身的命名（规范关系）。阶梯并非全错，而是范畴拼接：它把“能力的量”（第一级到第三级的纵轴）、“样本的常模”（人类在中间一级）、“造物的判据”（何谓专用）与“参照系的开放性”（何谓通用）焊在了一条线上。由此产生的“进步”叙事——系统沿阶梯上升——正是概念解绑论文所诊断的“进步转账错觉”——把工程进步记作智能增长——在层次图上的翻版\cite{yu2026twogoals}。

\subsection{每一类的“定义完成”形态}

类型论给出一个直接的推论表：能力概念可以操作定义收尾，但须附构造效度声明（否则退化为操作主义）；样本概念以常模收尾；造物概念以功能规格收尾；参照系概念\emph{永不}以定义收尾，只能以参照集声明收尾。四份“定义完成”的形态互不可替换——用造物规格去“完成”能力概念（第三种滑动）、用测量去“完成”本体（第一种滑动），都是类型错配。这张推论表同时解释了 AGI 定义竞争期（第~\ref{sec:survey} 节）的全部现象：竞争者各自拿一类概念的“定义完成”形态去回答另一类概念的问题。

\section{论题二：三层结构与三种非法滑动}\label{sec:layers}

\subsection{三层框架}

对每一个概念，区分三种主张：\textbf{本体层}（它是什么——存在与归属的主张）、\textbf{认识层}（什么证据支持“某系统具有它”——证据与推断的主张）、\textbf{方法层}（用什么程序检测它——操作与协议的主张）。三层不是同一件事的三种说法：它们有不同的真值条件、不同的失败方式、不同的修改代价。中文哲学话语对本体论—认识论—方法论三分有教科书级的稳定用法，本文直接沿用，并将每一处“$\times$ 在 $\times$ 层”的断言约束在该层内。

图~\ref{fig:matrix} 的矩阵是本文的工作图。定义之争的长期症结在图中表现为三种跨层滑动。

\subsection{滑动一：操作主义——测量冒充定义}

操作主义“概念即一套操作”的主张\cite{bridgman1927logic}在物理科学中已经退潮；这笔代价，心理测量学最先付出：测验确实测到了某个东西，但测验本身无法说明那是什么——Cronbach 与 Meehl 为此发明“构造效度”与“法则网络”来补救\cite{cronbach1955construct}。AI 的定义史中，这一教训三次重演。第一次：一个测量泛函（环境空间上的期望奖励）被当作智能的定义\cite{legg2007universal}。第二次：一套评测协议（技能获取效率、分级矩阵）被当作通用性的定义\cite{chollet2019measure,morris2023levels}。第三次：一份经济条款被当作 AGI 的定义\cite{openai2018charter}。

这里先作一处术语区分：分级框架标题中的“操作化”（operationalizing）\cite{morris2023levels}是评测工程的中性用法——把抽象目标翻译为可测指标；本文批评的操作主义是 Bridgman 意义上的——把测量操作当作概念的\emph{全部内容}。前者必要，后者越权；两者的分界正是认识层：评测工程停留在方法层并向认识层提交证据；操作主义跳过认识层，直接用方法层主张回答本体层问题。

操作主义定义的结构性缺陷在于它不可反驳。测试失败永远可以归因于方法——基准污染\cite{zhang2024gsm1k}、度量选择\cite{schaeffer2023mirage}、任务集过窄——而“定义”本身从不受伤。一个不承担失败风险的命题不是认识论主张。这也解释了引言中的那个直觉：“把方法论当认识论”的定义让人难以信服，因为信条不参与经验检验。

转机出现在 2021 至 2025 年：测量理论正以制度化的方式进入 AI 评测——基准的构造效度批评\cite{raji2021benchmark}、“涌现是度量伪影”的实证\cite{schaeffer2023mirage}、克隆基准暴露的过拟合\cite{zhang2024gsm1k}，直到 2025 年多篇论文把心理测量学的效度框架正式引入大模型评测\cite{bean2025measuring,salaudeen2025measurement,wallach2024evaluating,jo2025benchmark,xiao2023metriceval}。这条线为本文提供了最直接的支撑，但它全部停留在“评测—能力”链接上；经文献检索，未见把效度批评上推到“定义”层的系统诊断——这正是本文要填补的空白。哲学侧最接近的先行者是“AI 操作主义”的诊断\cite{mussgnug2024operationalism}——它论证机器学习实践把概念与操作化等同从而“暂时固定了指称”；本文与它方向相反：不是描述这一现象，而是给出防止它的分层规则与转译协议。

\subsection{滑动二：部分谬误——机制冒充能力定位}

第二种滑动发生在工程侧。Bennett 与 Hacker 的“部分谬误”（mereological fallacy）诊断：把整体层面的谓词归属于部分——“大脑思考”式的范畴错误\cite{bennett2003philosophical}。该诊断已被引入 AI 的“模型在思考”类语义归因讨论\cite{vantill2024personification,helliwell2024wittgenstein}，但尚未用于能力话语与评测分数的层次混淆——本文补上这一步。

论证分三步。第一，能力谓词的归属层级是系统与环境范围的组合（第~\ref{sec:types} 节）——“系统在环境范围 $E$ 内具有能力 $C$”的主语永远是系统，不是权重、电路或梯度。第二，通用性是关系属性\cite{yu2026theory}，而关系属性没有位置：说“泛化在权重中”不是错了，而是把一个不可定位的属性当成了可定位的东西——如同说“重量在公斤里”。第三，机制解释性研究合法产出的是\emph{机制归因}与\emph{实现证据}：找到某能力的实现机制（realizer），不等于给该能力找到了住所（location）——实现者携带倾向，倾向相对于任务分布成立，能力归属仍留在系统层。

为避免把转述的强化版本当作批评对象，本文以逐字引语建立被诊断对象的表述强度光谱（完整清单见附录~A）。纲领级：“特征是神经网络的基本单位”\cite{olah2020zoomin}——原文以思辨命题的形式提出，但它断言的是解释的基本单位在权重里，尚未断言能力在权重里。因果级：“归纳头是小规模纯注意力模型中上下文学习的主要机制”\cite{olsson2022induction}——限定范围内有消融证据支持的机制归因，合法；同文摘要的措辞则谨慎得多——“归纳头\emph{可能}构成该机制的初步且间接的证据”。实现级：“我们完整逆向工程了这些网络学到的算法”与三阶段中的“电路形成”\cite{nanda2023progress}——对玩具任务成立的最字面版本。许可级：“归因图给出的是关于底层模型中机制的\emph{假设}”“这些例子是存在性证明”\cite{lindsey2025biology}——作者自己把主张限制在认识层。描述级：“该概念宇宙在三个层次上有有趣的结构”\cite{tegmark2024geometry}——几何观察，未做概念断言。

光谱之外是\emph{漂移}。同一家机构的论文正文说“我们不能用归因图对底层模型得出确定性结论”\cite{lindsey2025biology}，面向公众的博客却说“我们发现 Claude 会提前规划”\cite{anthropic2025thoughts}——标题、正文、传播三级中，措辞强度逐级放大；社区转述再去掉限制词，把“可能是机制”转述成“就是机制”。实证侧已有反例：功能向量头而非归纳头主导较大模型的上下文学习\cite{yin2025which}；稀疏自编码器找不到规范的分析单元\cite{leask2025canonical}、跨层特征在版本间漂移而电路有效性随之下降\cite{anthropic2024crosscoders}；探针高分不等于表征中有结构\cite{hewitt2019designing,belinkov2022probing}；子空间补丁可制造“可解释性的错觉”\cite{makelov2023illusion}；因果探针方法在完备性与选择性间权衡\cite{canby2024reliable}。哲学侧，对机制解释的哲学研究已起步：有的论证链接不确定性比透明性更关键\cite{sullivan2022understanding}，有的区分透明的层次\cite{creel2022transparency}，有的规范机制发现的策略\cite{kastner2024explaining}，另有“机制解释性需要哲学”的倡言\cite{williams2025philosophy}与解释的数学哲学\cite{ayonrinde2025mathematical}，以及最接近本文的“因果主张必须披露识别假设”\cite{lin2026disclose}。但这些工作各覆盖一角——单点实证反证、探针限度、因果识别——没有人对“观察到结构$\to$能力在其中”这一\emph{推断形式}本身做分层诊断。
\subsection{滑动三：判据升格——造物判据冒充能力定义}

第三种滑动的后门在奠基文本中就已经打开：人工智能“制造优先”的定义结构（第~\ref{sec:survey} 节）把智能留作不加定义的初始词项。工程由此获得了一个正当的空间：不定义智能，只交付功能。但当造物的工程判据足够强势时——基准登顶、产品成功、经济回报——判据会溢出造物层，被当作能力概念乃至参照系概念的内容：把基准得分读作“智能水平”，把经济替代读作“通用性达成”\cite{openai2018charter}。与滑动一的区别在于载体：滑动一是“测量冒充定义”，滑动三是“工程目标冒充概念内容”；前者发生在概念的理论侧，后者发生在概念的应用侧。两者共享同一个缺口：跳过认识层——基准得分到能力主张之间缺构造效度声明\cite{raji2021benchmark,bean2025measuring}，经济替代到通用性主张之间缺实践抽样声明\cite{yu2026theory}。

\section{转译协议}\label{sec:protocol}

\subsection{为什么需要一份协议}

跨层转译本身并无不正当之处。工程观察必须能够进入概念层面的讨论——否则解释性研究与评测科学都失去意义；概念层面的讨论也必须能够约束工程——否则定义之争只是学院游戏。问题在于转译目前是\emph{隐含}发生的：最成熟的案例是 2025 年的 CHC 定义\cite{hendrycks2025definition}——它事实上完成了一次从心理测量模型到 AGI 概念的转译，但没有任何明确的转译方法论：读者看不到哪些心理学限制条件被带入、哪些被丢弃。转译协议的目的不是禁止转译，而是让每次转译\emph{出示条件}。它把本纲领方法论论文中的“术语对照”升格为一种方法论：那一篇做的是词与词的跨群体对照——把“采样轨迹（rollout）”对照为“实践流遥测”\cite{yu2026methodology}；本篇做的是命题与命题的跨层对照——工程观察陈述到概念命题，附生效条件与非法目录。

\subsection{转译表}

表~\ref{tab:protocol} 给出协议主体：六类常见工程观察，各自许可的概念命题与转译生效条件。

\begin{table}[t]
\centering
\caption{转译协议：工程观察$\to$概念命题的生效条件。}
\label{tab:protocol}
\small
\begin{tabularx}{\textwidth}{p{3.6cm} X p{4.2cm}}
\toprule
工程观察陈述 & 可转译成的概念命题 & 生效条件 \\
\midrule
消融电路 $c$ 后任务 $T$ 的表现下降 & $c$ 是系统在 $T$ 上的一个贡献机制（方法层主张） & 干预的完备性—选择性权衡已报告\cite{canby2024reliable}；因果识别假设已披露\cite{lin2026disclose} \\
探针从表征中解码出特征 $f$ & 表征中存在 $f$ 的可解码信息 & 控制任务选择性检验通过\cite{hewitt2019designing}；解码不等于被使用（遗忘式探测）\cite{belinkov2022probing} \\
特征空间出现几何结构（叶、晶体、方向） & 特征空间有描述性结构（描述级） & 结构的稳定性跨训练种子与字典尺寸检验\cite{leask2025canonical}；不得转译为“概念居住于此” \\
模型在基准 $B$ 上得分 $S$ & 相对 $B$ 的任务分布，系统有表现 $S$ & $B$ 附构造效度声明（测的是什么）与污染检查\cite{bean2025measuring,zhang2024gsm1k}；不得转译为“智能为 $S$” \\
系统在经济价值任务集 $W$ 上达到人类水平 & $W$ 上的替代证据（认识层） & $W$ 须明确声明为人类实践开放集的抽样\cite{yu2026theory}；抽样方式可审计\cite{yu2026methodology} \\
训练中出现相变（如“延迟泛化”，grokking\cite{power2022grokking}） & 训练动力学的一个阶段转变 & 相变的度量稳定性已检验\cite{schaeffer2023mirage}；不得自动转译为“泛化已在权重中形成”\cite{nanda2023progress} \\
\bottomrule
\end{tabularx}
\end{table}

协议的每行共享三条总则。其一，\emph{归属保持}：转译后的主语仍是系统（或系统与环境范围的组合），绝不变为权重、电路、梯度。其二，\emph{条件随行}：命题离开方法层时，它的生效条件一起走；无条件的概念命题是越权命题。其三，\emph{分层落位}：转译产出的命题须标注落位层——方法层主张（机制、信息）、认识层主张（证据、替代）、本体层主张（仅当条件全部满足且构造效度成立）。

\subsection{非法转译目录}

与转译表相对，四种反复出现的非法形式：\textbf{定位式}（“智能/泛化在 $X$ 中”）——给关系属性安位置，部分谬误；\textbf{整体错置式}（把系统层谓词归属于组件）——归属层级错误；\textbf{升格式}（“得分 $S$”$\to$“智能为 $S$”、“通过协议 $P$”$\to$“定义即 $P$”）——跳过构造效度，操作主义；\textbf{普遍化式}（从受测分布跳到开放集）——跳过抽样声明，把投影当原象。目录与转译表逐条相应：每一条非法形式都是某一条合法转译去掉生效条件后的产物。

\subsection{对 AGI 定义的推论：定义的三层写法}

协议对定义侧的推论：一份负责任的 AGI 定义应写成三层，而不是一句话。本体层：通用性是相对哪个实践集的关系属性（参照系声明）\cite{yu2026theory}；认识层：什么证据算“达到”——任务集、人群基准、时间窗口（认识声明）\cite{morris2023levels,hendrycks2025definition}；方法层：用什么程序测、抽样如何声明、失效条件是什么（方法声明）\cite{salaudeen2025measurement,yu2026methodology}。以此重读现有定义：分级矩阵是认识层与方法层上的扎实成果，但以“操作化”之名回答了本体层问题\cite{morris2023levels}；CHC 定义是本体层的一次认真声明，但未声明其样本投影的边界（CHC 是人类智能的样本结构，第~\ref{sec:types} 节）；经济定义是方法层的一纸合同，被当作了本体层裁决\cite{openai2018charter}。三层写法不消灭定义分歧——参照系之间的分歧没有实验裁决——但它让每一处分歧落到自己所属的层：可以点名、可以比较、可以撤回。

\section{在人本拟合论中的位置}\label{sec:program}

四词的最终图景在 AGI 人本拟合论中获得位置，也反过来支撑该理论。理论论文论证：AGI 是参照系概念（通用性相对于人类实践开放集）、其实现手段是对实践的拟合、其规范位置是增强性工具\cite{yu2026theory}。本文补全其概念底座：智能作为能力概念之所以是关系性的，是因为能力由拟合对象与实践范围共同定义——“发生上，这种智能靠拟合实践而获得”\cite{yu2026theory}；人类智能作为样本概念，是拟合论中人类参照系的测量面；人工智能作为造物概念，其判据（功能智能服务人类实践）已由概念解绑论文固定\cite{yu2026twogoals}——而该论文的“评估的不对称”（能力侧判据公开、机器之心一侧无共识测试）在本文的分层论中获得解释：能力谓词既有明确的系统层归属，又有现成的构造效度程序，而关于心的谓词缺少第三人称测试，故两轨分立。方法论论文的披露审计\cite{yu2026methodology}则示范了方法层的正确形态：不宣布测到了通用性，而是就其对八条流水线所披露的抽样方式与逐格证据作出声明——转译条件随行。

对四词的最终定位，一句话：\textbf{智能问度量，人类智能问常模，人工智能问造物判据，通用人工智能问参照系。}四个问题不能互相代替，正如四个概念不在一条阶梯上。

\section{证伪条款}\label{sec:falsify}

本文的三个主张各自可被驳倒。\textbf{类型论}：若未来工作证明四词可由单一标尺有意义地度量——一个度量同时对智能的能力结构、人类智能常模、人工智能的工程判据与 AGI 的参照集有不失真的预测力——则四类合并，类型论退化为一张多余的分类表；给出该度量的文献将取代本文。\textbf{分层论}：若发现某类能力谓词可以合法地归属于组件——即存在“组件级能力”的严格定义与不受反例动摇的判据——则部分谬误诊断失效，“泛化在权重中”从越权命题升格为待检验命题；机制解释性的进展若走到这一步，本文便成为被超越的先行者。\textbf{转译协议}：若四种非法转译形式之一被证明在可声明的条件下可靠（定位式在某种关系属性的可定位子类中成立，或升格式在某种自校准基准上成立），则目录相应行撤销、转译表相应扩充——协议自身附有可修订条款。空白声明同样可证伪：若存在中英文文献已做四词层次与边界的系统分析（以“概念类型”或等价框架），本文的增量贡献主张撤回，转为引用。

\section{结论}

四个词，四种概念，三层结构，一份协议。本文的诊断是：智能定义一个世纪不收敛、AGI 定义以换词回避裁决、工程话语把智能装进权重、公共话语难以取信于人——这四件事是同一种结构性缺陷的四种表现：概念主张在类型之间与层次之间未出示条件便滑动。类型论把四个词放回各自的座位；分层论给每个座位装上真值条件；转译协议在座位之间设了关卡——不是禁止通行，是查验证件。对工程侧，协议的产出是更诚实的强主张：机制归因、存在性证明、分布内倾向，每一件都值得做、值得说，只是不可冒充能力的住所。对定义侧，协议的产出是三层写法：参照系声明、认识声明、方法声明，各落各层，各赌各的证伪。对中文语境，本文补上一直缺少的那张图——“智力—智能—智慧”的直觉分层终于有了概念类型的学理形式。最终，四词在 AGI 人本拟合论中各就其位：智能是靠拟合实践而获得的能力，人类智能是参照系的测量面，人工智能是承载拟合的造物，通用人工智能是那个以人类实践开放集为准的、永不闭合的参照系本身。

\bibliographystyle{unsrtnat}
\bibliography{references}

\appendix

\section{主张清单：逐字引语与表述强度}\label{app:catalog}

本清单收录正文诊断所依据的逐字引语，按两个板块排列：定义侧（操作主义与判据升格的诊断对象）与机制侧（部分谬误的诊断对象）。“出处层级”标注引语所在的话语层级；“强度”为本文的分级（纲领级、因果级、实现级、许可级、描述级、传播级）；“许可”一栏给出：按第~\ref{sec:protocol} 节协议，该引语可支撑何种主张。

\subsection{定义侧}

\begin{longtable}{p{3.4cm} p{4.6cm} p{1.6cm} p{1.6cm} p{3.4cm}}
\toprule
出处 & 引语 & 出处层级 & 强度 & 许可 \\
\midrule
\endhead
Boring 1923\cite{boring1923intelligence} & ``Intelligence as the tests test it'' & 标题 & 操作主义原型 & 仅作历史诊断对象 \\
Legg \& Hutter 2007\cite{legg2007universal} & ``Intelligence measures an agent's ability to achieve goals in a wide range of environments'' & 正文 & 定义级 & 方法层泛函+条件，非本体定义 \\
Chollet 2019\cite{chollet2019measure} & 系统智能是技能获取效率的度量 & 摘要 & 定义级 & 方法层协议，附先验声明 \\
Morris 等 2023\cite{morris2023levels} & ``at least as capable as a human at most tasks'' & 正文 & 认识层框架 & 分级工件，不含本体裁决 \\
OpenAI 章程\cite{openai2018charter} & ``highly autonomous systems that outperform humans at most economically valuable work'' & 政策文本 & 方法层合同 & 经济替代证据，非通用性达成 \\
Hendrycks 等 2025\cite{hendrycks2025definition} & ``matching the cognitive versatility and proficiency of a well-educated adult'' & 摘要 & 本体层声明 & 参照系声明，缺样本边界声明 \\
Tesler 定理\cite{tesler1997effect} & ``AI is whatever hasn't been done yet'' & 短文 & 否定性类别 & 仅作移动边界的历史注脚 \\
\bottomrule
\end{longtable}

\subsection{机制侧}

\begin{longtable}{p{3.2cm} p{4.8cm} p{1.7cm} p{1.5cm} p{3.0cm}}
\toprule
出处 & 引语 & 出处层级 & 强度 & 许可 \\
\midrule
\endhead
Olah 等 2020\cite{olah2020zoomin} & ``Features are the fundamental unit of neural networks'' & 论文正文 & 纲领级 & 解释单位的纲领，未及能力定位 \\
Olsson 等 2022（摘要）\cite{olsson2022induction} & ``preliminary and indirect evidence for a hypothesis that induction heads might constitute the mechanism'' & 摘要 & 许可级 & 假设级证据声明（合规样本） \\
Olsson 等 2022（正文）\cite{olsson2022induction} & 归纳头是小规模纯注意力模型中上下文学习的主要机制 & 正文（限定范围） & 因果级 & 限定范围内机制归因 \\
Nanda 等 2023\cite{nanda2023progress} & ``We fully reverse engineer the algorithm learned by these networks''；三阶段中的“电路形成” & 摘要 & 实现级 & 玩具任务的实现证据 \\
Varma 等 2023\cite{varma2023circuit} & “延迟泛化”（grokking）源于记忆电路与泛化电路的竞争 & 摘要 & 因果级 & 合成设置的机制解释 \\
Wang 等 2024\cite{wang2024grokked} & ``the formation of the generalizing circuit'' & 摘要 & 实现级 & 训练动力学主张 \\
Lindsey 等 2025（正文）\cite{lindsey2025biology} & ``attribution graphs provide hypotheses about mechanisms''；``existence proofs'' & 正文 & 许可级 & 认识层的自我设限（合规样本） \\
Anthropic 博客 2025\cite{anthropic2025thoughts} & ``we found that Claude plans ahead'' & 机构博客 & 传播级 & 传播层强主张（诊断对象） \\
Tegmark 团队 2025\cite{tegmark2024geometry} & ``this concept universe has interesting structure at three levels'' & 摘要 & 描述级 & 几何描述，不作概念断言 \\
Park 等 2024\cite{park2024linear} & ``high-level concepts are represented linearly as directions'' & 摘要（非正式表述） & 假说级 & 形式化假说+证据 \\
Gurnee \& Tegmark 2024\cite{gurnee2024space} & 标题断言``Represent''，摘要降为``suggest $\ldots$ basic ingredients of a world model'' & 标题/摘要 & 标题强、摘要弱 & 漂移样本 \\
Yin \& Steinhardt 2025\cite{yin2025which} & 上下文学习主要依赖功能向量头而非归纳头 & 摘要 & 实证反证 & 对强读法的反例 \\
Leask 等 2025\cite{leask2025canonical} & 稀疏自编码器找不到规范的分析单元 & 标题 & 实证反证 & 特征非规范性的证据 \\
\bottomrule
\end{longtable}

清单的读法：机制侧前八行构成“纲领$\to$因果$\to$实现$\to$许可$\to$传播”的强度光谱，后五行是描述级、假说级、标题—正文漂移样本与实证反证。

\section{中文语境对接}\label{app:cn}

\subsection{术语史时间线}

1922--1924：比纳量表经修订进入中国，心理测量固定“智力”译名；1978--1981：“人工智能”译名在大陆学界启用，中国人工智能学会筹建；2002--2003：北京大学设中国第一个智能科学系与“智能科学与技术”专业；2022：教育部学科目录将“智能科学与技术”列为交叉学科门类一级学科。“智力”（心理学）与“智能”（智能科学）的分工沿用至今。

\subsection{政策用词演变}

政策用词演变的细节见第~\ref{sec:cn} 节，此处记录要点：2017 年规划在基础理论布局中提及“通用人工智能”\cite{state2017plan}；2023 年北京政策把通用人工智能等同于通用大模型技术体系\cite{beijing2023measures}；以 2025 年政府工作报告为代表的中央文件改用“人工智能+”“具身智能”“智能经济”等重点用词，未使用“通用人工智能”\cite{gov2025report}。政策语言的漂移与国际上的换词漂移（前沿模型、超级智能）同构：操作词汇绕开定义之争。

\subsection{对话者定位}

朱松纯团队的“无限任务、自主定义任务、价值驱动”三特征与通智测试\cite{zhu2025tongtest,zhu2023origin,zhu2025moment}是中文语境最强的 AGI 定义体系——按本文的分层，属本体层声明与认识层工件；刘永谋的“唯名论困境”\cite{liu2024complex}是中文哲学侧对定义缺失的元批评，本文的类型论是对其“困境”的结构性回应（困境源于要求四类概念共用一种“定义完成”形态）；张钹的三代论\cite{zhang2020third}与谭铁牛的官方定义句式\cite{tan2024lecture}分别提供历史分期与政策基准；钟义信的机制主义\cite{zhong2018mechanism}给出信息论进路的智能定义，其“本体论信息/认识论信息”术语与本文三层框架相互呼应；赵汀阳\cite{zhao2019selfconscious}从主体性维度提出的“新主体”问题，在本文框架中落位于参照系概念的规范面。空白声明：以四词层次与边界为题的系统分析，中文文献未见；若存在，本文即以之为准并撤回空白主张。

\end{document}
